%% 編譯器：XeLaTeX（Overleaf: Menu -> Compiler -> XeLaTeX）
\documentclass[12pt]{article}
\usepackage[UTF8]{ctex}
\usepackage[margin=1in]{geometry}
\usepackage{booktabs,array,longtable,hyperref,enumitem}
\usepackage{xurl}
\onehalfspacing
\newcommand{\path}[1]{\texttt{\detokenize{#1}}}

\title{機器學習與金融科技 --- 期末專題}
\author{}
\date{}

\begin{document}
\maketitle

\section{複製論文 (Replicating a High-Quality Paper)}

Feng, Fuli, Xiangnan He, Xiang Wang, Cheng Luo, Yiqun Liu, and Tat-Seng Chua.
``Temporal Relational Ranking for Stock Prediction.''
\textit{ACM Transactions on Information Systems} 37, no.~2 (2019): Article 27.
\url{https://doi.org/10.1145/3309547}.

\begin{itemize}[leftmargin=*]
  \item 期刊等級：交大資工 A 級期刊清單（\path{共用參考資料} 或課程 GitHub
        \path{journal-ranking/交大資工A級期刊_20200910Updated.xlsx} 第 11 筆）
  \item 官方程式碼：\url{https://github.com/fulifeng/Temporal_Relational_Stock_Ranking}
  \item 本機複製環境：\path{../Temporal_Relational_Stock_Ranking}（已 clone、已修補 TF1$\to$TF2 相容性、
        venv 已建立、baseline RankLSTM 已實測跑通）
  \item \textbf{跑訓練指令前務必加 \path{TF_USE_LEGACY_KERAS=1}}，例如
        \path{TF_USE_LEGACY_KERAS=1 python rank_lstm.py -p ../data/2013-01-01 -m NASDAQ -l 4 -u 32}，
        否則會在 \path{BasicLSTMCell} 那行出現
        \path{AttributeError: BasicLSTMCell is not available with Keras 3}。
        細節見 \path{報告/R3_模型與實驗設計/PROGRESS.md}。
\end{itemize}

\section{論文摘要與方法}

提出 Relational Stock Ranking (RSR) 模型與 Temporal Graph Convolution，將股票間的關係
（產業關係、Wikidata 關係）以「時間敏感」的方式編碼進排序任務，用於預測股票隔日報酬排名，
取代過去把股票視為互相獨立個體的做法。

\section{真正的評分依據（以 syllabus 與官方 GitHub 為準）}

\path{報告/R1~R4} 底下原本寫「對應評分表項目」，該評分表查無出處：兩份 syllabus
（\path{slides/20260907-Ch00-Syllabus.pdf}、\path{slides/20260914-Ch00-Syllabus.pdf}）與官方 GitHub
（\path{202609-ML-FinTech/00-course-info}，2026-09-14 以 GitHub API 查證）皆無此表，
也沒有任何「R1/R2/R3/R4 分節評分」的說法。以下為查證過的依據。

\paragraph{Grading policy}（\path{20260907-Ch00-Syllabus.pdf} p.6）

\begin{center}
\begin{tabular}{lrp{8cm}}
\toprule
項目 & \% & 內容 \\
\midrule
Participation & 20 & 課堂練習、簡報（project/MFS/HW）、課堂總結，會 cold call \\
\textbf{Project} & \textbf{40} & \textbf{根據 manuscript（LaTeX 撰寫）評分}，另需準備簡報用的 slides \\
Exam & 40 & 課堂考，可帶一張 A4 雙面小抄 \\
\bottomrule
\end{tabular}
\end{center}

\paragraph{個人 repo 結構}（同上 p.5）

\begin{verbatim}
202609-ML-FinTech-<studentID>-<nickname>/
|-- HW/
`-- Project/   <- README.md, logs, snapshots, data and codes
\end{verbatim}

\paragraph{Project 的 README.md 要包含}（p.10 ``Your GitHub Repo''）：
GitHub 頁面連結、Chicago style 引用複製的論文、Overleaf manuscript 連結、Canva slides 連結、
``Code''（Jupyter Notebook）、``Data''（資料集或連結）。

真正被評分的是 Project 40\% 那份 LaTeX manuscript，不是這個資料夾本身。
\path{報告/R1~R4} 的內容（動機、EDA、模型實驗、結論）仍是 manuscript 需要的材料，
只是拿掉「各自佔幾 \% 評分」的錯誤框架，現在當成寫 manuscript 前的草稿分節：

\begin{center}
\begin{longtable}{>{\raggedright\arraybackslash}p{4.6cm}>{\raggedright\arraybackslash}p{4cm}>{\raggedright\arraybackslash}p{6cm}}
\toprule
資料夾 & manuscript 對應章節 & 內容 \\
\midrule
\endhead
\path{報告/R1_主題與動機} & Introduction / Motivation & 論文介紹、為何選這篇、RQ1/RQ2 \\
\path{報告/R2_資料與EDA} & Data & 原始資料說明、EDA notebook（\path{eda.ipynb}） \\
\path{報告/R3_模型與實驗設計} & Methods / Experiments & Baseline (Rank\_LSTM) vs.\ RSR，含 ablation（有無關係圖、動態延伸） \\
\path{報告/R4_實證分析與結論} & Results / Conclusion & 結果解讀、與原論文比較、延伸討論 \\
\path{文獻_論文參考} & References & 複製論文全文、相關文獻筆記 \\
\path{程式碼} & --- & 指向 \path{../Temporal_Relational_Stock_Ranking}，或放整理過的 Jupyter Notebook 版本 \\
\path{資料} & --- & 資料集或資料連結說明 \\
\bottomrule
\end{longtable}
\end{center}

\section{延伸方向（已確認）：動態跨股票關聯}

RSR 論文的 industry / wiki 關係圖在訓練期間固定不變；HGTAN（\path{../HGTAN}）的群體 hypergraph
歸屬（產業、基金持股）也不隨時間更新。確認要做的延伸：在 \path{報告/R3_模型與實驗設計} 的 ablation
額外加入一組「動態關聯」設定，用滾動窗口相關係數（或其他隨時間重估的關係圖）取代固定關係圖，
重跑 \path{relation_rank_lstm.py} 比較排名表現，驗證「關聯的變化本身」是否比「關聯的存在與否」更有預測力。

對照課程 \path{20260914-Ch00-Syllabus.pdf} 「Potentials of your projects」投影片，
這屬於 ``with modifications (new dataset or new methods)''，是可以再往 thesis / journal publication 走的路線。

\section{資料}

\begin{itemize}[leftmargin=*]
  \item \textbf{Sequential Data}：NASDAQ/NYSE 逾 8,000 檔股票 30 年歷史日線資料
        （open, high, low, close, volume），來源 Google Finance，論文使用 2013-01-01 起的處理後版本
  \item \textbf{Industry Relation}：NASDAQ/NYSE 股票的產業關係（sector/industry）
  \item \textbf{Wiki Relation}：從 Wikidata 抽取的公司關係
  \item 資料已隨官方 repo 提供，路徑：\path{../Temporal_Relational_Stock_Ranking/data/}
\end{itemize}

\section{Code}

Python，Jupyter Notebook（README 要求交 Jupyter Notebook，原始訓練腳本為 \path{.py}，另外整理成 \path{.ipynb}）：
\begin{itemize}[leftmargin=*]
  \item \path{training/rank_lstm.py}：Baseline，Rank\_LSTM（不含關係圖）
  \item \path{training/relation_rank_lstm.py}：完整模型，Relational Stock Ranking (RSR)
\end{itemize}

\section{Manuscript / Slides / GitHub}

\begin{itemize}[leftmargin=*]
  \item Overleaf：（待補連結，分享至 venteng@gmail.com）
  \item Canva Slides：（待補連結，分享至 venteng@gmail.com）
  \item GitHub：（待補：學生自己的 GitHub 頁面連結）
\end{itemize}

\end{document}
